\documentclass[a4paper,12pt]{letter}  % tipo do documento
\usepackage[brazil]{babel}           % define a linguagem padrão
\usepackage{graphicx}                % permite a inserção de figuras
\usepackage[T1]{fontenc}
% \usepackage[latin1]{inputenc}        % para usar caracteres acentuados diretamente
\usepackage[utf8]{inputenc}
\usepackage[hmargin=2cm,vmargin=2cm]{geometry}
\usepackage{listings} %para colocar codigos
\usepackage{amsmath}
\usepackage{amssymb}
\lstset{numbers=left,
stepnumber=1,
firstnumber=1,
numberstyle=\tiny,
extendedchars=true,
breaklines=true,
frame=tb,
basicstyle=\footnotesize,
% stringstyle=\ttfamily,
% showstringspaces=false
}




% \newcommand{\gabarito}[1]{\textbf{\begin{small} \\ #1 \\ \end{small}}} %mostrar gabarito
\newcommand{\gabarito}[1]{\textbf{\begin{small} \end{small}}} %nao mostrar gabarito


\begin{document}
UNIVERSIDADE FEDERAL DO CEARÁ\\
Departamento Estatística e Matemática Aplicada\\
Disciplina de Elementos de Análise Combinatória\\
Professor: Albert Einstein F. Muritiba.\\

%Temas: Princípios de Contagem, Permutações, Combinações, Triângulo de Pascal, Binômio e Polinômios\\[0.5cm]
% Nome do Aluno: \makebox[280pt]{\hrulefill} \vspace{2pt}
%Matr\'icula:\makebox[60pt]{\hrulefill} \vspace{2pt}\\
\begin{center}
	Lista de exercícios
	%   2ª Avaliação - Segunda Chamada
	% //npc 2 %+ nef
	% NEF
\end{center}

\begin{enumerate}
	\item Faça como no exemplo:
	\begin{enumerate}
        \item $\sum_{i=1}^5 i^2 = 1^2 +2^2 +3^2 +4^2 +5^2$ (exemplo)
        \item $\sum_{i=1}^5 2.i = $
        \item $\sum_{i=1}^5 2^i = $
        \item $\sum_{i=1}^5 (x+i)^3 = $
        \item $3 + 6+ 9 + 12 + 15 = \sum_{i=1}^5 3.i$ (exemplo)
        \item $ 1/1+ 1/2+1/3+1/4+1/5 = $
        \item $ 12 + 13 + 14 +15 + 16 = $
        \item $ (x+1) + 2(x+2) + 3(x+3) + 4(x+4) +5(x+5) =$ 
        \item $ 1 - 2 + 3 - 4 + 5 =$ 
        \item $\sum_{i=1}^5 i^3 = \sum_{i=1}^4 i^3 + 5^3$ (exemplo)
        \item $\sum_{i=1}^5 5.i =$ 
        \item $\sum_{i=1}^{10} (x+1)^i =$ 
    \end{enumerate}
	\item para todo natural $n\ge 1$ prove por indução que $P(n)$ é verdade:
	\begin{enumerate}
        \item P(n): $\sum_{i=1}^n (3i-1) = \frac{n(3n+1)}{2}$
        \item P(n): $\sum_{i=1}^n 2(3i-2) = 3n^2-n$
        \item P(n): $\sum_{i=1}^n \frac{1}{i(i+1)} = \frac{n}{n+1}$
        \item P(n): $\sum_{i=1}^n 2^n = 2(2^n-1)$
        \item P(n): $\sum_{i=1}^n i.i! = (n+1)!-1$
        \item P(n): $4^n -1$ é divisível por 3
        \item P(n): $2^{2n-1} +1$ é divisível por 3
        \item P(n): $17n^3 + 103n$ é divisível por 6
        % \item P(n): $2^n \le (n+1)!$ 
        % \item P(n): $3^n \le (n+2)!$ 
        % \item P(n): $2^n \ge n+1$ 
        \item P(n): $\sum_{i=1}^n F_{2i-1} = F_{2n}$
        \item P(n): $\sum_{i=1}^n F_{2i} = F_{2n+2}-1$
    \end{enumerate}
	\item 

\end{enumerate}

\end{document}
